\documentclass[11pt]{article}
\usepackage{amsmath, amssymb, physics, mathtools, bm}
\usepackage{tikz, pgfplots}
\pgfplotsset{compat=1.18}
\usepackage[margin=2.3cm]{geometry}
\usepackage{siunitx, booktabs, hyperref}
\usepackage{braket}

\newcommand{\D}{\mathrm{D}}
\newcommand{\de}{\mathrm{d}}
\newcommand{\pd}[2]{\frac{\partial #1}{\partial #2}}
\newcommand{\pdd}[2]{\frac{\partial^{2} #1}{\partial #2^{2}}}
\newcommand{\Lap}{\nabla^{2}}

\title{B5 General Relativity \& Cosmology --- Model Solutions, 2019}
\author{}
\date{}

\begin{document}
\maketitle

%======================================================================
\section*{Question 1 --- Equivalence principle and geodesic equation}

\subsection*{(a) Equivalence principle and proper time}

\textbf{Problem.} \emph{State the equivalence principle. In a free-falling lab with coordinates $\xi^\alpha$, a particle obeys $\de^2\xi^\alpha/\de\tau^2 = 0$. Define the proper time $\tau$.}

\textbf{Equivalence principle (strong form):} at every spacetime event one can choose a local inertial (free-falling) frame in which the laws of physics take their special-relativistic form, and gravitational effects are locally indistinguishable from acceleration.

Free-falling frame: $\eta_{\alpha\beta}=\mathrm{diag}(-1,+1,+1,+1)$. Invariant line element:
\[
\de s^2 = \eta_{\alpha\beta}\,\de\xi^\alpha\,\de\xi^\beta = -(\de\xi^0)^2 + (\de\bm\xi)^2.
\]
For a timelike worldline, $\de s^2 < 0$. Define proper time as the interval read by a clock comoving with the particle:
\[
c^2\,\de\tau^2 \equiv -\de s^2 = -\eta_{\alpha\beta}\,\de\xi^\alpha\,\de\xi^\beta.
\]
Divide by $c^2$:
\[
\boxed{\;\de\tau^2 = -\tfrac{1}{c^2}\,\eta_{\alpha\beta}\,\de\xi^\alpha\,\de\xi^\beta.\;}
\]

\subsection*{(b) Transformation to general coordinates}

\textbf{Problem.} \emph{Starting from $\de^2\xi^\alpha/\de\tau^2 = 0$, derive the geodesic equation in coordinates $x^\alpha$ and identify the symmetry of $\Gamma$.}

Chain rule on $\xi^\alpha = \xi^\alpha(x^\mu)$:
\[
\frac{\de \xi^\alpha}{\de\tau} = \pd{\xi^\alpha}{x^\mu}\,\frac{\de x^\mu}{\de\tau}.
\]
Differentiate w.r.t.\ $\tau$, applying product rule:
\[
\frac{\de^2\xi^\alpha}{\de\tau^2} = \frac{\de}{\de\tau}\!\left(\pd{\xi^\alpha}{x^\mu}\right)\!\frac{\de x^\mu}{\de\tau} + \pd{\xi^\alpha}{x^\mu}\,\frac{\de^2 x^\mu}{\de\tau^2}.
\]
Chain rule on $\partial\xi^\alpha/\partial x^\mu$ along the curve:
\[
\frac{\de}{\de\tau}\!\left(\pd{\xi^\alpha}{x^\mu}\right) = \frac{\partial^2\xi^\alpha}{\partial x^\nu\partial x^\mu}\,\frac{\de x^\nu}{\de\tau}.
\]
Combine and set equal to zero:
\[
\pd{\xi^\alpha}{x^\mu}\,\frac{\de^2 x^\mu}{\de\tau^2} + \frac{\partial^2\xi^\alpha}{\partial x^\mu\partial x^\nu}\,\frac{\de x^\mu}{\de\tau}\frac{\de x^\nu}{\de\tau}=0.
\]
Multiply by $\partial x^\lambda/\partial\xi^\alpha$ (sum on $\alpha$):
\[
\pd{x^\lambda}{\xi^\alpha}\pd{\xi^\alpha}{x^\mu}\,\frac{\de^2 x^\mu}{\de\tau^2} + \pd{x^\lambda}{\xi^\alpha}\frac{\partial^2\xi^\alpha}{\partial x^\mu\partial x^\nu}\,\frac{\de x^\mu}{\de\tau}\frac{\de x^\nu}{\de\tau} = 0.
\]
Apply $(\partial x^\lambda/\partial\xi^\alpha)(\partial\xi^\alpha/\partial x^\mu)=\delta^\lambda_\mu$:
\[
\delta^\lambda_\mu\,\frac{\de^2 x^\mu}{\de\tau^2} + \pd{x^\lambda}{\xi^\alpha}\frac{\partial^2\xi^\alpha}{\partial x^\mu\partial x^\nu}\,\frac{\de x^\mu}{\de\tau}\frac{\de x^\nu}{\de\tau} = 0.
\]
Contract the Kronecker delta:
\[
\frac{\de^2 x^\lambda}{\de\tau^2} + \pd{x^\lambda}{\xi^\alpha}\frac{\partial^2\xi^\alpha}{\partial x^\mu\partial x^\nu}\,\frac{\de x^\mu}{\de\tau}\frac{\de x^\nu}{\de\tau} = 0.
\]
Identify the connection:
\[
\Gamma^\lambda_{\mu\nu} \equiv \pd{x^\lambda}{\xi^\alpha}\frac{\partial^2\xi^\alpha}{\partial x^\mu\partial x^\nu}.
\]
Geodesic equation:
\[
\boxed{\;\frac{\de^2 x^\lambda}{\de\tau^2} + \Gamma^\lambda_{\mu\nu}\,\frac{\de x^\mu}{\de\tau}\frac{\de x^\nu}{\de\tau} = 0.\;}
\]
Partial derivatives commute, hence $\Gamma^\lambda_{\mu\nu} = \Gamma^\lambda_{\nu\mu}$ (symmetric in lower indices).

\subsection*{(c) Metric compatibility relation}

\textbf{Problem.} \emph{With $g_{\mu\nu}=\eta_{\alpha\beta}(\partial\xi^\alpha/\partial x^\mu)(\partial\xi^\beta/\partial x^\nu)$, show $\partial_\lambda g_{\mu\nu} = g_{\rho\nu}\Gamma^\rho_{\lambda\mu} + g_{\mu\rho}\Gamma^\rho_{\lambda\nu}$.}

Invariant line element:
\[
\de s^2 = \eta_{\alpha\beta}\,\de\xi^\alpha\,\de\xi^\beta = g_{\mu\nu}\,\de x^\mu\,\de x^\nu,\qquad g_{\mu\nu} = \eta_{\alpha\beta}\,\pd{\xi^\alpha}{x^\mu}\pd{\xi^\beta}{x^\nu}.
\]
Differentiate $g_{\mu\nu}$ w.r.t.\ $x^\lambda$, using product rule:
\[
\pd{g_{\mu\nu}}{x^\lambda} = \eta_{\alpha\beta}\,\frac{\partial^2\xi^\alpha}{\partial x^\lambda\partial x^\mu}\pd{\xi^\beta}{x^\nu} + \eta_{\alpha\beta}\,\pd{\xi^\alpha}{x^\mu}\frac{\partial^2\xi^\beta}{\partial x^\lambda\partial x^\nu}.
\]
Invert the connection definition: multiply $\Gamma^\lambda_{\mu\nu}=(\partial x^\lambda/\partial\xi^\alpha)(\partial^2\xi^\alpha/\partial x^\mu\partial x^\nu)$ by $\partial\xi^\beta/\partial x^\lambda$ and apply $(\partial\xi^\beta/\partial x^\lambda)(\partial x^\lambda/\partial\xi^\alpha) = \delta^\beta_\alpha$:
\[
\frac{\partial^2\xi^\alpha}{\partial x^\mu\partial x^\nu} = \Gamma^\rho_{\mu\nu}\,\pd{\xi^\alpha}{x^\rho}.
\]
Substitute this identity (with index labels $(\lambda,\mu)$ and $(\lambda,\nu)$):
\[
\pd{g_{\mu\nu}}{x^\lambda} = \eta_{\alpha\beta}\,\Gamma^\rho_{\lambda\mu}\pd{\xi^\alpha}{x^\rho}\pd{\xi^\beta}{x^\nu} + \eta_{\alpha\beta}\,\pd{\xi^\alpha}{x^\mu}\,\Gamma^\rho_{\lambda\nu}\pd{\xi^\beta}{x^\rho}.
\]
Recognise $g_{\rho\nu}$ and $g_{\mu\rho}$ from the metric definition:
\[
\boxed{\;\pd{g_{\mu\nu}}{x^\lambda} = g_{\rho\nu}\,\Gamma^\rho_{\lambda\mu} + g_{\mu\rho}\,\Gamma^\rho_{\lambda\nu}.\;}
\]
\label{eq:metric-compat}

\subsection*{(d) Christoffel formula}

\textbf{Problem.} \emph{Use the symmetry of $\Gamma$ to express $\Gamma^\sigma_{\mu\lambda}$ in terms of $g_{\mu\nu}$.}

Restate \eqref{eq:metric-compat}:
\[
\partial_\lambda g_{\mu\nu} = g_{\rho\nu}\Gamma^\rho_{\lambda\mu} + g_{\mu\rho}\Gamma^\rho_{\lambda\nu}.
\]
Cyclic permutation $\lambda\to\mu\to\nu\to\lambda$:
\[
\partial_\mu g_{\nu\lambda} = g_{\rho\lambda}\Gamma^\rho_{\mu\nu} + g_{\nu\rho}\Gamma^\rho_{\mu\lambda}.
\]
Cyclic permutation once more:
\[
\partial_\nu g_{\lambda\mu} = g_{\rho\mu}\Gamma^\rho_{\nu\lambda} + g_{\lambda\rho}\Gamma^\rho_{\nu\mu}.
\]
Form $\partial_\mu g_{\nu\lambda} + \partial_\nu g_{\lambda\mu} - \partial_\lambda g_{\mu\nu}$:
\[
\partial_\mu g_{\nu\lambda} + \partial_\nu g_{\lambda\mu} - \partial_\lambda g_{\mu\nu} = g_{\rho\lambda}\Gamma^\rho_{\mu\nu} + g_{\nu\rho}\Gamma^\rho_{\mu\lambda} + g_{\rho\mu}\Gamma^\rho_{\nu\lambda} + g_{\lambda\rho}\Gamma^\rho_{\nu\mu} - g_{\rho\nu}\Gamma^\rho_{\lambda\mu} - g_{\mu\rho}\Gamma^\rho_{\lambda\nu}.
\]
Apply $\Gamma^\rho_{\mu\nu}=\Gamma^\rho_{\nu\mu}$ to cancel the mixed pairs:
\[
g_{\nu\rho}\Gamma^\rho_{\mu\lambda} - g_{\mu\rho}\Gamma^\rho_{\lambda\nu} = 0,\qquad g_{\rho\mu}\Gamma^\rho_{\nu\lambda} - g_{\rho\nu}\Gamma^\rho_{\lambda\mu} = 0.
\]
Surviving terms (both equal $g_{\rho\lambda}\Gamma^\rho_{\mu\nu}$ by symmetry):
\[
\partial_\mu g_{\nu\lambda} + \partial_\nu g_{\lambda\mu} - \partial_\lambda g_{\mu\nu} = 2\,g_{\rho\lambda}\,\Gamma^\rho_{\mu\nu}.
\]
Multiply by $\tfrac{1}{2}g^{\sigma\lambda}$ and use $g^{\sigma\lambda}g_{\rho\lambda}=\delta^\sigma_\rho$:
\[
\tfrac{1}{2}g^{\sigma\lambda}\!\left(\partial_\mu g_{\nu\lambda} + \partial_\nu g_{\lambda\mu} - \partial_\lambda g_{\mu\nu}\right) = \Gamma^\sigma_{\mu\nu}.
\]
Relabel $\nu\to\lambda$, $\lambda\to\nu$ to match the requested form:
\[
\boxed{\;\Gamma^\sigma_{\mu\lambda} = \tfrac{1}{2}\,g^{\nu\sigma}\!\left(\pd{g_{\mu\nu}}{x^\lambda} + \pd{g_{\lambda\nu}}{x^\mu} - \pd{g_{\lambda\mu}}{x^\nu}\right).\;}
\]

\subsection*{(e) Covariant derivative}

\textbf{Problem.} \emph{Define $U^\lambda = \de x^\lambda/\de\tau$. Use the geodesic equation to define the covariant derivative; comment on its rank and significance.}

Geodesic equation rewritten using $U^\lambda$:
\[
\frac{\de U^\lambda}{\de\tau} + \Gamma^\lambda_{\mu\nu}\,U^\mu U^\nu = 0.
\]
Apply chain rule along the curve: $\de U^\lambda/\de\tau = (\partial_\nu U^\lambda)(\de x^\nu/\de\tau) = U^\nu\partial_\nu U^\lambda$:
\[
U^\nu\,\partial_\nu U^\lambda + \Gamma^\lambda_{\mu\nu}\,U^\mu U^\nu = 0.
\]
Factor $U^\nu$:
\[
U^\nu\!\left(\partial_\nu U^\lambda + \Gamma^\lambda_{\mu\nu}\,U^\mu\right) = 0.
\]
Identify the bracket as the covariant derivative:
\[
\boxed{\;\nabla_\nu U^\lambda \equiv \partial_\nu U^\lambda + \Gamma^\lambda_{\mu\nu}\,U^\mu.\;}
\]
\textbf{Significance.} $\nabla_\nu U^\lambda$ corrects the partial derivative for the rotation/stretching of the coordinate basis induced by curvature, so that geodesic motion ($\nabla_U U = 0$) is the curved-space analogue of Newton's first law: a free particle parallel-transports its own four-velocity.

\textbf{Tensor character.} Under $x^\mu\to x'^\mu$, transform $\partial_\nu U^\lambda$ using product rule on the Jacobian:
\[
\partial_\nu U^\lambda \;\to\; \pd{x^\beta}{x'^\nu}\!\left(\pd{x'^\lambda}{x^\alpha}\partial_\beta U^\alpha + U^\alpha\frac{\partial^2 x'^\lambda}{\partial x^\alpha\partial x^\beta}\right).
\]
Transform $\Gamma^\lambda_{\mu\nu}U^\mu$ (the inhomogeneous transformation law of $\Gamma$ contributes a second-derivative piece with opposite sign):
\[
\Gamma^\lambda_{\mu\nu}U^\mu \;\to\; \pd{x'^\lambda}{x^\alpha}\pd{x^\beta}{x'^\nu}\Gamma^\alpha_{\gamma\beta}U^\gamma - \pd{x^\beta}{x'^\nu}U^\alpha\frac{\partial^2 x'^\lambda}{\partial x^\alpha\partial x^\beta}.
\]
Add: the non-tensorial second-derivative pieces cancel:
\[
\nabla_\nu U^\lambda \;\to\; \pd{x'^\lambda}{x^\alpha}\pd{x^\beta}{x'^\nu}\,\nabla_\beta U^\alpha.
\]
Hence $\nabla_\nu U^\lambda$ transforms as a $(1,1)$ mixed rank-2 tensor (one upper Jacobian, one lower inverse Jacobian) --- required so that $\nabla_U U = 0$ has covariant meaning.

%======================================================================
\section*{Question 2 --- Schwarzschild geodesics and orbits}

\subsection*{(a) Geodesic equations from the Lagrangian}

\textbf{Problem.} \emph{Write the geodesic equations for the Schwarzschild line element using the Lagrange procedure; explain the $\theta$ equation; relate $p^\alpha$ to the geodesics.}

Effective Lagrangian (one-half line element):
\[
\mathcal L = \tfrac{1}{2}g_{\mu\nu}\dot x^\mu\dot x^\nu = \tfrac{1}{2}\!\left[-\!\left(1-\tfrac{2GM}{c^2 r}\right)c^2\dot t^2 + \!\left(1-\tfrac{2GM}{c^2 r}\right)^{-1}\!\!\dot r^2 + r^2\dot\theta^2 + r^2\sin^2\theta\,\dot\phi^2\right].
\]
Euler--Lagrange equations $\frac{\de}{\de\tau}\partial_{\dot x^\mu}\mathcal L - \partial_{x^\mu}\mathcal L = 0$:
\[
\frac{\de}{\de\tau}\!\left[\!\left(1-\tfrac{2GM}{c^2 r}\right)c^2\dot t\right] = 0,
\]
\[
\frac{\de}{\de\tau}\!\left[\!\left(1-\tfrac{2GM}{c^2 r}\right)^{-1}\!\dot r\right] - \tfrac{1}{2}\!\left[-\tfrac{2GM}{c^2 r^2}c^2\dot t^2 - \tfrac{2GM}{c^2 r^2}\!\left(1-\tfrac{2GM}{c^2 r}\right)^{-2}\!\dot r^2 + 2r\dot\theta^2 + 2r\sin^2\theta\,\dot\phi^2\right] = 0,
\]
\[
\frac{\de}{\de\tau}\!\left[r^2\dot\theta\right] - r^2\sin\theta\cos\theta\,\dot\phi^2 = 0,
\]
\[
\frac{\de}{\de\tau}\!\left[r^2\sin^2\theta\,\dot\phi\right] = 0.
\]
Spherical symmetry: orient axes so initial conditions are $\theta = \pi/2$, $\dot\theta = 0$. Insert into the $\theta$-equation:
\[
\frac{\de}{\de\tau}(r^2\cdot 0) - r^2\sin(\pi/2)\cos(\pi/2)\dot\phi^2 = 0.
\]
Use $\cos(\pi/2)=0$:
\[
\ddot\theta = 0.
\]
Hence $\theta(\tau) = \pi/2$ is preserved identically and the orbit lies in the equatorial plane.

\textbf{Four-momentum.} For a massive particle $p^\alpha = m\,\de x^\alpha/\de\tau$; multiplying the geodesic equation by $m$ gives the parallel-transport law $p^\nu\nabla_\nu p^\alpha = 0$, i.e.\ free particles follow geodesic curves with their four-momentum tangent to the geodesic.

\subsection*{(b) Affine parameter and constants of motion}

\textbf{Problem.} \emph{When is $\tau$ an appropriate affine parameter? Identify the conserved quantities $k$ and $h$ from the $t$ and $\phi$ equations.}

Proper time $\tau$ is affine only for \emph{timelike} (massive) geodesics, where $\de\tau^2 = -\de s^2/c^2 > 0$. For \emph{null} (massless) geodesics $\de\tau = 0$, so pick any affine parameter $\lambda$ along which $p^\alpha = \de x^\alpha/\de\lambda$ and the geodesic equation $\de p^\alpha/\de\lambda + \Gamma^\alpha_{\mu\nu}p^\mu(\de x^\nu/\de\lambda) = 0$ holds.

Integrate the $t$-equation (Lagrangian independent of $t$, so $\partial_{\dot t}\mathcal L$ conserved):
\[
\!\left(1-\tfrac{2GM}{c^2 r}\right)c^2\dot t = \text{const} \equiv c^2 k.
\]
Integrate the $\phi$-equation in the equatorial plane ($\sin^2\theta = 1$):
\[
r^2\dot\phi = \text{const} \equiv h.
\]
Lower the index using $p_\alpha = g_{\alpha\beta}p^\beta = m g_{\alpha\beta}\dot x^\beta$. For $\alpha=0$ use $g_{tt}=-(1-2GM/c^2r)c^2$:
\[
p_0 = -m\!\left(1-\tfrac{2GM}{c^2 r}\right)c^2\dot t = -mc^2 k.
\]
For $\alpha=3$ use $g_{\phi\phi}=r^2\sin^2\theta$:
\[
p_3 = m r^2\sin^2\theta\,\dot\phi = m h.
\]
\textbf{Physical meaning.} $k = -p_0/(mc^2)$ is the conserved \emph{energy per unit rest mass} measured at infinity (associated with Killing vector $\partial_t$ via time-translation symmetry); $h = p_3/m$ is the conserved \emph{specific angular momentum} (Killing vector $\partial_\phi$ via rotational symmetry).

\subsection*{(c) Energy equation}

\textbf{Problem.} \emph{From $g_{\mu\nu}\dot x^\mu\dot x^\nu = c^2$ derive the radial energy equation.}

Divide $-c^2\de\tau^2 = g_{\mu\nu}\de x^\mu\de x^\nu$ by $\de\tau^2$, set $\dot\theta=0$, $\theta=\pi/2$:
\[
-\!\left(1-\tfrac{2GM}{c^2 r}\right)c^2\dot t^2 + \!\left(1-\tfrac{2GM}{c^2 r}\right)^{-1}\!\dot r^2 + r^2\dot\phi^2 = -c^2.
\]
Solve for $\dot t$ from $k=(1-2GM/c^2r)\dot t$:
\[
\dot t = \frac{k}{1-2GM/(c^2 r)}.
\]
Substitute $\dot t$ and $\dot\phi = h/r^2$:
\[
-\!\left(1-\tfrac{2GM}{c^2 r}\right)\frac{c^2 k^2}{(1-2GM/c^2 r)^2} + \!\left(1-\tfrac{2GM}{c^2 r}\right)^{-1}\!\dot r^2 + \frac{h^2}{r^2} = -c^2.
\]
Simplify the first term:
\[
-\frac{c^2 k^2}{1-2GM/(c^2 r)} + \frac{\dot r^2}{1-2GM/(c^2 r)} + \frac{h^2}{r^2} = -c^2.
\]
Multiply through by $(1-2GM/(c^2 r))$:
\[
-c^2 k^2 + \dot r^2 + \frac{h^2}{r^2}\!\left(1-\tfrac{2GM}{c^2 r}\right) = -c^2\!\left(1-\tfrac{2GM}{c^2 r}\right).
\]
Expand the right-hand side:
\[
-c^2 k^2 + \dot r^2 + \frac{h^2}{r^2}\!\left(1-\tfrac{2GM}{c^2 r}\right) = -c^2 + \frac{2GM}{r}.
\]
Rearrange (move $-c^2 k^2$ and $2GM/r$ across):
\[
\boxed{\;\dot r^2 + \frac{h^2}{r^2}\!\left(1-\tfrac{2GM}{c^2 r}\right) - \frac{2GM}{r} = c^2(k^2 - 1).\;}
\]
\label{eq:energy}

\subsection*{(d) Orbit equation and relativistic correction}

\textbf{Problem.} \emph{Using \eqref{eq:energy} and $r^2\dot\phi=h$ derive $\de^2 u/\de\phi^2 + u = GM/h^2 + 3GM u^2/c^2$ with $u=1/r$.}

Change variable $u = 1/r$, so $r = 1/u$ and $\de r/\de\phi = -u^{-2}\de u/\de\phi$. Use chain rule with $\dot\phi = h/r^2 = hu^2$:
\[
\dot r = \frac{\de r}{\de\phi}\,\dot\phi = -\frac{1}{u^2}\frac{\de u}{\de\phi}\cdot h u^2.
\]
Cancel $u^2$:
\[
\dot r = -h\,\frac{\de u}{\de\phi}.
\]
Square:
\[
\dot r^2 = h^2\!\left(\frac{\de u}{\de\phi}\right)^{\!2}.
\]
Substitute $\dot r^2$ and $1/r = u$ into \eqref{eq:energy}:
\[
h^2\!\left(\frac{\de u}{\de\phi}\right)^{\!2} + h^2 u^2\!\left(1-\tfrac{2GM u}{c^2}\right) - 2GM u = c^2(k^2-1).
\]
Expand the bracket on the second term:
\[
h^2\!\left(\frac{\de u}{\de\phi}\right)^{\!2} + h^2 u^2 - \frac{2GM h^2}{c^2}u^3 - 2GM u = c^2(k^2-1).
\]
Differentiate w.r.t.\ $\phi$ (RHS is constant):
\[
2h^2\,\frac{\de u}{\de\phi}\frac{\de^2 u}{\de\phi^2} + 2h^2 u\,\frac{\de u}{\de\phi} - \frac{6GM h^2}{c^2}u^2\,\frac{\de u}{\de\phi} - 2GM\frac{\de u}{\de\phi} = 0.
\]
Factor out $2\,\de u/\de\phi$:
\[
2\,\frac{\de u}{\de\phi}\!\left[h^2\frac{\de^2 u}{\de\phi^2} + h^2 u - \frac{3GM h^2}{c^2}u^2 - GM\right] = 0.
\]
Non-circular orbit: $\de u/\de\phi \neq 0$ so the bracket vanishes:
\[
h^2\frac{\de^2 u}{\de\phi^2} + h^2 u - \frac{3GM h^2}{c^2}u^2 - GM = 0.
\]
Divide by $h^2$:
\[
\frac{\de^2 u}{\de\phi^2} + u - \frac{3GM}{c^2}u^2 - \frac{GM}{h^2} = 0.
\]
Rearrange to the requested form:
\[
\boxed{\;\frac{\de^2 u}{\de\phi^2} + u = \frac{GM}{h^2} + \frac{3GM}{c^2}u^2.\;}
\]
\label{eq:orbit}
The Newtonian equation has only the constant source $GM/h^2$. The relativistic correction is the term $3GM u^2/c^2$. Its ratio to the Newtonian term is $3h^2 u^2/c^2 = 3 v_\perp^2/c^2$ where $v_\perp = h u$, so it is significant when $v\sim c$, i.e.\ for tight orbits at $r\sim r_s = 2GM/c^2$ (e.g.\ Mercury's perihelion advance, S2 around Sgr A$^*$).

\subsection*{(e) Effective potential in GR}

\textbf{Problem.} \emph{Cast the energy equation in Newtonian form $\tfrac{1}{2}(\de r/\de t)^2 + V_{\mathrm{eff}}(r) = E$.}

In the orbiting observer's frame the affine parameter is $\tau$, so $\de r/\de t$ is replaced by $\de r/\de\tau = \dot r$. Multiply \eqref{eq:energy} by $\tfrac{1}{2}$:
\[
\tfrac{1}{2}\dot r^2 + \tfrac{1}{2}\!\left[\frac{h^2}{r^2}\!\left(1-\tfrac{2GM}{c^2 r}\right) - \frac{2GM}{r}\right] = \tfrac{1}{2}c^2(k^2-1).
\]
Expand the inner bracket:
\[
\tfrac{1}{2}\dot r^2 + \frac{h^2}{2r^2} - \frac{GM h^2}{c^2 r^3} - \frac{GM}{r} = \tfrac{1}{2}c^2(k^2-1).
\]
Match to the form $\tfrac{1}{2}\dot r^2 + V_{\mathrm{eff}}(r) = E$:
\[
\boxed{\;V_{\mathrm{eff}}(r) = -\frac{GM}{r} + \frac{h^2}{2r^2} - \frac{GM h^2}{c^2 r^3},\qquad E = \tfrac{1}{2}c^2(k^2-1).\;}
\]
\label{eq:Veff}
The first two terms reproduce Newtonian gravity + centrifugal barrier; the $1/r^3$ term is purely relativistic.

\subsection*{(f) Circular orbits}

\textbf{Problem.} \emph{Show circular-orbit extrema occur at $r = (h/2GM)\!\left(h\pm\sqrt{h^2 - 12 G^2M^2/c^2}\right)$.}

Circular orbit condition: $\de V_{\mathrm{eff}}/\de r = 0$. Differentiate \eqref{eq:Veff} term by term:
\[
\frac{\de V_{\mathrm{eff}}}{\de r} = \frac{GM}{r^2} - \frac{h^2}{r^3} + \frac{3GM h^2}{c^2 r^4}.
\]
Set to zero:
\[
\frac{GM}{r^2} - \frac{h^2}{r^3} + \frac{3GM h^2}{c^2 r^4} = 0.
\]
Multiply by $r^4$:
\[
GM\,r^2 - h^2\,r + \frac{3GM h^2}{c^2} = 0.
\]
Solve the quadratic in $r$ via the standard formula:
\[
r = \frac{h^2 \pm \sqrt{h^4 - 4\cdot GM\cdot 3GM h^2/c^2}}{2GM}.
\]
Simplify under the surd:
\[
r = \frac{h^2 \pm \sqrt{h^4 - 12\,G^2 M^2 h^2/c^2}}{2GM}.
\]
Factor $h^2$ inside the surd:
\[
r = \frac{h^2 \pm h\sqrt{h^2 - 12\,G^2 M^2/c^2}}{2GM}.
\]
Pull out the common factor $h$ in the numerator:
\[
\boxed{\;r = \frac{h}{2GM}\!\left(h \pm \sqrt{h^2 - \frac{12\,G^2M^2}{c^2}}\right).\;}
\]
The $+$ root is the stable minimum of $V_{\mathrm{eff}}$ (regular orbit), the $-$ root is the unstable maximum (light-ring / unstable circular orbit). Real solutions require $h^2 \geq 12 G^2 M^2/c^2$; equality gives the innermost stable circular orbit (ISCO) at $r = 6GM/c^2$.

%======================================================================
\section*{Question 3 --- Binary pulsar and gravitational radiation}

\subsection*{(a) Why a GR expert should notice}

\textbf{Problem.} \emph{Pulsar in a binary, $T = \SI{0.10225156248}{day}$. Why does this matter for GR?}

A binary system containing a pulsar provides a clean clock embedded in a strong-field, dynamical spacetime: the orbit decays via emission of gravitational waves, and the highly accurate pulse timing lets one measure $\dot T$ at the level needed to test the quadrupole formula of GR (Hulse--Taylor / double pulsar). It is one of very few systems testing the radiative sector of Einstein's theory.

\subsection*{(b) Orbital radius}

\textbf{Problem.} \emph{$T_0 = \SI{0.10225156248}{day}$, $\dot T \approx -10^{-9}\ \mathrm{day}/\SI{1000}{day}$, $m_1 = m_2 = 1.4 M_\odot$, circular orbit. Find $a$.}

Convert period to SI:
\[
T = 0.10225156248\times 86400\,\mathrm{s} = \SI{8834.535}{s}.
\]
Apply Kepler's third law (total mass $M_{\mathrm{tot}} = 2 m_{\mathrm p} = 2.8\,M_\odot$, $G M_\odot = \SI{1.327e20}{m^3 s^{-2}}$):
\[
a^3 = \frac{G M_{\mathrm{tot}} T^2}{4\pi^2}.
\]
Evaluate $G M_{\mathrm{tot}}$:
\[
G M_{\mathrm{tot}} = 2.8\times \SI{1.327e20}{m^3 s^{-2}} = \SI{3.716e20}{m^3 s^{-2}}.
\]
Evaluate $T^2$:
\[
T^2 = (\SI{8834.535}{s})^2 = \SI{7.805e7}{s^2}.
\]
Multiply:
\[
G M_{\mathrm{tot}} T^2 = \SI{3.716e20}{m^3 s^{-2}}\times\SI{7.805e7}{s^2} = \SI{2.901e28}{m^3}.
\]
Divide by $4\pi^2$:
\[
a^3 = \frac{\SI{2.901e28}{m^3}}{4\pi^2} = \SI{7.346e26}{m^3}.
\]
Take the cube root:
\[
a = (\SI{7.346e26}{m^3})^{1/3} \approx \SI{9.02e8}{m}.
\]
Compare to $R_\odot = \SI{6.96e8}{m}$:
\[
\boxed{\;a \approx \SI{9.02e8}{m} \approx 1.3\,R_\odot.\;}
\]

\subsection*{(c) Definitions in the quadrupole formula}

\textbf{Problem.} \emph{Define $\bar h^{ij}$ and $t'$ in $\bar h^{ij} = (2G/c^6 r)\,\de^2 I^{ij}/\de t'^2$.}

$\bar h^{\mu\nu}$ is the \emph{trace-reversed} metric perturbation in the harmonic gauge:
\[
g_{\mu\nu} = \eta_{\mu\nu} + h_{\mu\nu},\qquad \bar h_{\mu\nu} = h_{\mu\nu} - \tfrac{1}{2}\eta_{\mu\nu}\,h,\qquad h \equiv \eta^{\mu\nu}h_{\mu\nu}.
\]
The harmonic-gauge condition $\partial_\mu\bar h^{\mu\nu} = 0$ reduces the linearised Einstein equation to $\Box\bar h^{\mu\nu} = (16\pi G/c^4)T^{\mu\nu}$ (sign consistent with the paper's $-8\pi G/c^4$ convention for the Einstein equation), whose retarded solution gives the quadrupole formula.

$t'$ is the \emph{retarded time} at the source:
\[
\boxed{\;t' = t - r/c,\;}
\]
the time of emission for a wave received at the field point $(t,\bm r)$.

\subsection*{(d) Meaning of $x^i, x^j$}

\textbf{Problem.} \emph{What do $x^i, x^j$ in $I^{ij} = M c^2 x^i x^j$ represent, given that pulsars are point-like?}

Pulsars are $\sim\!\SI{10}{km}$ across against orbital separations of $\sim\!\SI{e9}{m}$, so each behaves as a point mass and the internal structure does not contribute to $I^{ij}$. For two equal masses $M_{\mathrm p} = M/2$ on a circular orbit, each pulsar sits at $\pm\bm x(t')$ with $|\bm x| = a/2$ on opposite sides of the centre of mass. The standard quadrupole sum is
\[
I^{ij} \;=\; \sum_{A=1,2} M_{\mathrm p} c^2\,x_A^i x_A^j \;=\; 2 M_{\mathrm p}c^2\,x^i x^j \;=\; M c^2\,x^i x^j.
\]
Hence $x^i$ are the \emph{Cartesian components of one pulsar's orbital position vector about the centre of mass}, of magnitude $a/2$. They label the macroscopic orbital trajectory, not internal pulsar coordinates.

\subsection*{(e) GW frequency and waveform}

\textbf{Problem.} \emph{Find $f_{\mathrm{GW}}$ to 2 s.f.; write $\bar h^{ij}(t')$.}

Orbital angular frequency:
\[
\omega_{\mathrm{orb}} = \frac{2\pi}{T} = \frac{2\pi}{\SI{8834.535}{s}} = \SI{7.111e-4}{rad/s}.
\]
Orbital frequency:
\[
f_{\mathrm{orb}} = \frac{1}{T} = \frac{1}{\SI{8834.535}{s}} = \SI{1.132e-4}{Hz}.
\]
Quadrupole radiation: $f_{\mathrm{GW}} = 2 f_{\mathrm{orb}}$:
\[
f_{\mathrm{GW}} = 2\times \SI{1.132e-4}{Hz} = \SI{2.26e-4}{Hz}.
\]
Round to 2 s.f.:
\[
\boxed{\;f_{\mathrm{GW}} \approx \SI{2.3e-4}{Hz}.\;}
\]
\textbf{Waveform.} Place orbit in $xy$-plane. Pulsar 1 sits at $\bm x = (a/2)(\cos\omega t', \sin\omega t', 0)$; pulsar 2 at $-\bm x$. Sum quadrupole contributions:
\[
I^{ij} = M_{\mathrm p}c^2\,x^i x^j + M_{\mathrm p}c^2\,(-x)^i(-x)^j = 2 M_{\mathrm p}c^2\,x^i x^j = M c^2\,x^i x^j,
\]
matching the formula in the question (with $M = 2 M_{\mathrm p}$). Take $\omega \equiv \omega_{\mathrm{orb}}$ and compute components:
\[
x^1(t') = (a/2)\cos(\omega t'),\qquad x^2(t') = (a/2)\sin(\omega t'),\qquad x^3 = 0.
\]
Hence:
\[
I^{11} = M c^2 (a/2)^2\cos^2(\omega t'),
\]
\[
I^{22} = M c^2 (a/2)^2\sin^2(\omega t'),
\]
\[
I^{12} = M c^2 (a/2)^2\sin(\omega t')\cos(\omega t').
\]
Apply $\cos^2\theta = \tfrac12(1+\cos 2\theta)$:
\[
I^{11} = \tfrac{1}{2}M c^2 (a/2)^2\,[1+\cos(2\omega t')].
\]
Apply $\sin^2\theta = \tfrac12(1-\cos 2\theta)$:
\[
I^{22} = \tfrac{1}{2}M c^2 (a/2)^2\,[1-\cos(2\omega t')].
\]
Apply $\sin\theta\cos\theta = \tfrac12\sin 2\theta$:
\[
I^{12} = \tfrac{1}{2}M c^2 (a/2)^2\sin(2\omega t').
\]
Differentiate twice w.r.t.\ $t'$ (constants drop; $\de^2/\de t'^2[\cos 2\omega t']=-4\omega^2\cos 2\omega t'$):
\[
\ddot I^{11} = -\tfrac{1}{2}M c^2 (a/2)^2 \cdot 4\omega^2\cos(2\omega t').
\]
Collect $(a/2)^2\cdot 4 = a^2$:
\[
\ddot I^{11} = -\tfrac{1}{2}M c^2 a^2 \omega^2\cos(2\omega t').
\]
Similarly $\ddot I^{22}$ (sign flips because $1-\cos 2\omega t'$):
\[
\ddot I^{22} = +\tfrac{1}{2}M c^2 a^2\omega^2\cos(2\omega t').
\]
And $\ddot I^{12}$ (second derivative of $\sin 2\omega t'$ is $-4\omega^2\sin 2\omega t'$):
\[
\ddot I^{12} = -\tfrac{1}{2}M c^2 a^2\omega^2\sin(2\omega t').
\]
Insert into $\bar h^{ij} = (2G/c^6 r)\ddot I^{ij}$:
\[
\bar h^{11} = \frac{2G}{c^6 r}\cdot\!\left(-\tfrac{1}{2}M c^2 a^2 \omega^2\cos 2\omega t'\right).
\]
Cancel the factor of $2$ and reduce $c^2/c^6 = 1/c^4$:
\[
\bar h^{11} = -\frac{G M a^2\omega^2}{c^4 r}\cos(2\omega t').
\]
Similarly for the remaining components:
\[
\bar h^{22} = +\frac{G M a^2\omega^2}{c^4 r}\cos(2\omega t'),\qquad \bar h^{12} = \bar h^{21} = -\frac{G M a^2\omega^2}{c^4 r}\sin(2\omega t').
\]
Collect into a matrix:
\[
\boxed{\;\bar h^{ij}(t') = \frac{G\,M\,a^2\,\omega^2}{c^4\,r}\begin{pmatrix}-\cos 2\omega t' & -\sin 2\omega t' & 0\\ -\sin 2\omega t' & +\cos 2\omega t' & 0 \\ 0 & 0 & 0\end{pmatrix},\;}
\]
oscillating at angular frequency $2\omega_{\mathrm{orb}}$, hence frequency $\SI{2.3e-4}{Hz}$.

\subsection*{(f) Why $\dot T$ tests GR}

\textbf{Problem.} \emph{Why does measuring $\dot T$ test general relativity?}

Energy radiated as gravitational waves, summed over the two polarisations and integrated over the sphere, is the Einstein quadrupole luminosity formula (with reduced quadrupole $Q^{ij} = M x^i x^j -\tfrac13\delta^{ij}\delta_{kl}M x^k x^l$, i.e.\ $I^{ij}/c^2$ minus trace):
\[
P_{\mathrm{GW}} = \frac{G}{5 c^5}\,\langle\dddot Q_{ij}\dddot Q^{ij}\rangle.
\]
For a circular binary of masses $m_1,m_2$ and separation $a$ this evaluates to
\[
P_{\mathrm{GW}} = \frac{32}{5}\frac{G^4}{c^5}\frac{(m_1 m_2)^2(m_1+m_2)}{a^5}.
\]
This luminosity drains the orbital binding energy $E_{\mathrm{orb}} = -G m_1 m_2/(2a)$. Energy balance $\dot E_{\mathrm{orb}} = -P_{\mathrm{GW}}$ shrinks $a$:
\[
\dot a = -\frac{64}{5}\frac{G^3}{c^5}\frac{m_1 m_2(m_1+m_2)}{a^3}.
\]
Kepler's third law $T^2 \propto a^3/(m_1+m_2)$ gives $\dot T/T = \tfrac{3}{2}\dot a/a$, so
\[
\dot T = -\frac{192\pi}{5 c^5}\!\left(\frac{2\pi G}{T}\right)^{\!5/3}\!\frac{m_1 m_2}{(m_1+m_2)^{1/3}}.
\]
$\dot T$ is therefore predicted with \emph{no free parameters} once $m_1,m_2$ are known. The slight eccentricity permits independent measurement of post-Keplerian parameters (periastron advance, Shapiro delay, gravitational redshift) which over-determine $m_1, m_2$. Comparing predicted $\dot T$ to the measured one is thus a direct, parameter-free test of GR's radiative quadrupole sector --- famously passed by Hulse--Taylor and by the double pulsar to better than $0.1\%$.

%======================================================================
\section*{Question 4 --- Cosmological constant and the FRW universe}

\subsection*{(a) Bianchi identities}

\textbf{Problem.} \emph{State the contracted Bianchi identities and verify that the modified Einstein equation respects them. Why is this physically significant?}

Contracted Bianchi identity (Einstein tensor is divergence-free):
\[
\nabla^\mu\!\left(R_{\mu\nu} - \tfrac{1}{2}g_{\mu\nu}R\right) = 0.
\]
Take the covariant divergence $\nabla^\mu$ of equation (2):
\[
\nabla^\mu\!\left(R_{\mu\nu} - \tfrac{1}{2}g_{\mu\nu}R\right) + \Lambda\,\nabla^\mu g_{\mu\nu} = -\frac{8\pi G}{c^4}\nabla^\mu T_{\mu\nu}.
\]
Apply Bianchi (first term vanishes) and metric compatibility $\nabla^\mu g_{\mu\nu}=0$ (second term vanishes):
\[
0 + 0 = -\frac{8\pi G}{c^4}\,\nabla^\mu T_{\mu\nu}.
\]
Conclude:
\[
\nabla^\mu T_{\mu\nu} = 0.
\]
\textbf{Significance.} The modified equation enforces local conservation of energy--momentum automatically. Adding $\Lambda g_{\mu\nu}$ does not spoil this because $g_{\mu\nu}$ is covariantly constant; any other proposed addition would have to share that property.

\subsection*{(b) Four-velocity of the cosmic fluid}

\textbf{Problem.} \emph{Define $U^\mu$ in $T_{\mu\nu} = P g_{\mu\nu}+(\rho + P/c^2)U_\mu U_\nu$. Give $U^\lambda, U_\lambda$ at $\lambda=t,r$ in the FRW comoving frame.}

$U^\mu = \de x^\mu/\de\tau$ is the four-velocity of a fluid element measured in its own proper time, normalised $g_{\mu\nu}U^\mu U^\nu = -c^2$.

Comoving frame: fluid at rest in spatial coordinates, $\de r = \de\theta = \de\phi = 0$. Reduce line element:
\[
-c^2\de\tau^2 = -c^2\de t^2.
\]
Solve for $\de t/\de\tau$ (take positive root):
\[
\frac{\de t}{\de\tau} = 1.
\]
Identify the four-velocity components:
\[
U^t = \frac{\de t}{\de\tau} = 1,\qquad U^r = \frac{\de r}{\de\tau} = 0.
\]
Read off the diagonal metric components: $g_{tt} = -c^2$, $g_{rr} = R^2(t)/(1-r^2/a^2)$:
\[
U_t = g_{tt}U^t = (-c^2)(1) = -c^2.
\]
\[
U_r = g_{rr}U^r = \frac{R^2}{1-r^2/a^2}\cdot 0 = 0.
\]
Collect:
\[
\boxed{\;U^t = 1,\ U^r = 0;\qquad U_t = -c^2,\ U_r = 0.\;}
\]

\subsection*{(c) Effective fluid with $\Lambda$}

\textbf{Problem.} \emph{Show that with $\rho\to\rho', P\to P'$ the modified equation reduces to the original; evaluate $\rho_V$ and interpret.}

Form the primed stress tensor with $\rho\to\rho'$, $P\to P'$:
\[
T'_{\mu\nu} = P'\,g_{\mu\nu}+(\rho'+P'/c^2)U_\mu U_\nu.
\]
Substitute $\rho' = \rho + \rho_V$ and $P' = P - \rho_V c^2$:
\[
T'_{\mu\nu} = (P-\rho_V c^2)\,g_{\mu\nu} + \bigl[(\rho + \rho_V) + (P-\rho_V c^2)/c^2\bigr]U_\mu U_\nu.
\]
Simplify the bracket: the $+\rho_V$ and $-\rho_V c^2/c^2$ cancel:
\[
(\rho + \rho_V) + (P-\rho_V c^2)/c^2 = \rho + P/c^2.
\]
Hence:
\[
T'_{\mu\nu} = (P-\rho_V c^2)\,g_{\mu\nu} + (\rho + P/c^2)U_\mu U_\nu.
\]
Split using $T_{\mu\nu} = P g_{\mu\nu} + (\rho+P/c^2)U_\mu U_\nu$:
\[
T'_{\mu\nu} = T_{\mu\nu} - \rho_V c^2\,g_{\mu\nu}.
\]
Insert into equation (1) with $T_{\mu\nu}\to T'_{\mu\nu}$:
\[
R_{\mu\nu}-\tfrac{1}{2}g_{\mu\nu}R = -\frac{8\pi G}{c^4}T'_{\mu\nu}.
\]
Expand the RHS:
\[
R_{\mu\nu}-\tfrac{1}{2}g_{\mu\nu}R = -\frac{8\pi G}{c^4}T_{\mu\nu} + \frac{8\pi G}{c^2}\rho_V\,g_{\mu\nu}.
\]
Rearrange the modified equation (2) into the same form:
\[
R_{\mu\nu}-\tfrac{1}{2}g_{\mu\nu}R = -\frac{8\pi G}{c^4}T_{\mu\nu} - \Lambda\,g_{\mu\nu}.
\]
Match the coefficient of $g_{\mu\nu}$:
\[
\frac{8\pi G}{c^2}\rho_V = -\Lambda.
\]
Solve for $\rho_V$:
\[
\boxed{\;\rho_V = -\frac{\Lambda c^2}{8\pi G}.\;}
\]

\textbf{Interpretation.} $\rho_V$ is the energy density of the \emph{vacuum}: a Lorentz-invariant perfect fluid with equation of state $P_V = -\rho_V c^2$, constant in space and time, indistinguishable from the cosmological-constant term in the field equations.

\subsection*{(d) FRW identity $\ddot R/R = H^2 + \dot H$}

\textbf{Problem.} \emph{Prove $\ddot R/R = H^2 + \dot H$.}

Definition $H = \dot R/R$ implies $\dot R = H R$. Differentiate w.r.t.\ $t$:
\[
\ddot R = \dot H R + H\dot R.
\]
Substitute $\dot R = H R$:
\[
\ddot R = \dot H R + H^2 R.
\]
Divide by $R$:
\[
\boxed{\;\frac{\ddot R}{R} = H^2 + \dot H.\;}
\]
\label{eq:Hidentity}

\subsection*{(e) Differential equation for $H$}

\textbf{Problem.} \emph{For a flat $\Lambda$+matter universe show $2\dot H + 3H^2 = \Lambda c^2$.}

For non-relativistic matter, $p = 0$. The two Friedmann equations reduce to
\[
\frac{\ddot R}{R} = -\frac{4\pi G}{3}\rho + \tfrac{1}{3}\Lambda c^2,
\]
\[
H^2 = \frac{8\pi G}{3}\rho + \tfrac{1}{3}\Lambda c^2 \quad (k=0).
\]
Eliminate $\rho$. Multiply the second by $\tfrac{1}{2}$:
\[
\tfrac{1}{2}H^2 = \frac{4\pi G}{3}\rho + \tfrac{1}{6}\Lambda c^2.
\]
Add to the first:
\[
\frac{\ddot R}{R} + \tfrac{1}{2}H^2 = \tfrac{1}{2}\Lambda c^2.
\]
Use \eqref{eq:Hidentity}:
\[
\dot H + H^2 + \tfrac{1}{2}H^2 = \tfrac{1}{2}\Lambda c^2.
\]
Multiply by $2$:
\[
\boxed{\;2\dot H + 3H^2 = \Lambda c^2.\;}
\]
\label{eq:Hode}

\subsection*{(f) Solving for $H(t)$}

\textbf{Problem.} \emph{If $H\propto 1/t$ at early times, show $H = A\coth(3At/2)$ with $A^2 = \Lambda c^2/3$, and comment on late-time behaviour.}

Early-time check (matter domination, $\Lambda$ negligible): \eqref{eq:Hode} becomes $2\dot H + 3H^2 = 0$. Trial $H = c_1/t$ gives $-2c_1/t^2 + 3c_1^2/t^2 = 0$, hence $c_1 = 2/3$:
\[
H_{\mathrm{early}} = \frac{2}{3t} \propto \frac{1}{t}.
\]
This matches the hypothesis and fixes the boundary condition $H\to\infty$ as $t\to 0^+$.

Define $A^2 \equiv \Lambda c^2/3$ and rearrange \eqref{eq:Hode}:
\[
2\,\frac{\de H}{\de t} = \Lambda c^2 - 3 H^2.
\]
Factor the RHS:
\[
2\,\frac{\de H}{\de t} = 3(A^2 - H^2).
\]
Separate variables:
\[
\frac{2\,\de H}{A^2 - H^2} = 3\,\de t.
\]
Integrate using the hint $\int\de x/(p^2-x^2) = (1/p)\coth^{-1}(x/p)$ with $p=A$, $x=H$:
\[
\frac{2}{A}\coth^{-1}(H/A) = 3 t + \mathrm{const}.
\]
Apply $H\to\infty$ as $t\to 0^+$, so $\coth^{-1}(H/A)\to 0$ and the constant is $0$:
\[
\frac{2}{A}\coth^{-1}(H/A) = 3 t.
\]
Solve for $\coth^{-1}(H/A)$:
\[
\coth^{-1}(H/A) = \tfrac{3At}{2}.
\]
Invert:
\[
\boxed{\;H(t) = A\coth(3At/2),\qquad A = \sqrt{\Lambda c^2/3}.\;}
\]
\textbf{Late time.} $\coth(3At/2)\to 1$ as $t\to\infty$:
\[
H \to A = \sqrt{\Lambda c^2/3}.
\]
The Hubble parameter approaches a constant, the universe enters de Sitter (exponential) expansion $R\propto e^{At}$, dominated by the cosmological constant.

\subsection*{(g) Big rip from phantom vacuum $\rho_V\propto R^\epsilon$}

\textbf{Problem.} \emph{If $\rho_V\propto R^\epsilon$ with $\epsilon > 0$, show that $R(t)$ diverges in finite time in a flat universe (matter neglected after $t_0$).}

Flat universe, ordinary matter neglected after $t_0$. Friedmann equation reduces to
\[
H^2 = \frac{8\pi G}{3}\rho_V.
\]
Substitute $\rho_V = \rho_{V,0}(R/R_0)^\epsilon$:
\[
H^2 = \frac{8\pi G}{3}\rho_{V,0}\,(R/R_0)^\epsilon.
\]
Define $B^2 \equiv (8\pi G/3)\rho_{V,0}\,R_0^{-\epsilon}$:
\[
H^2 = B^2\,R^\epsilon.
\]
Take positive root (expanding universe, $H>0$):
\[
H = B\,R^{\epsilon/2}.
\]
Substitute $H = \dot R/R$:
\[
\frac{\dot R}{R} = B\,R^{\epsilon/2}.
\]
Rearrange (multiply by $\de t$, divide by $R^{\epsilon/2}$):
\[
R^{-1-\epsilon/2}\,\de R = B\,\de t.
\]
Integrate from $(t_0, R_0)$ to $(t, R)$:
\[
\int_{R_0}^{R} R'^{-1-\epsilon/2}\,\de R' = \int_{t_0}^{t}B\,\de t'.
\]
Evaluate (power rule, $\epsilon>0$):
\[
\left[\frac{R'^{-\epsilon/2}}{-\epsilon/2}\right]_{R_0}^{R} = B(t-t_0).
\]
Simplify the prefactor:
\[
-\frac{2}{\epsilon}\!\left[R^{-\epsilon/2} - R_0^{-\epsilon/2}\right] = B(t - t_0).
\]
Solve for $R^{-\epsilon/2}$:
\[
R^{-\epsilon/2} = R_0^{-\epsilon/2} - \tfrac{1}{2}\epsilon B(t - t_0).
\]
$R\to\infty$ corresponds to $R^{-\epsilon/2}\to 0$:
\[
0 = R_0^{-\epsilon/2} - \tfrac{1}{2}\epsilon B(t_{\mathrm{rip}}-t_0).
\]
Solve for $t_{\mathrm{rip}}-t_0$:
\[
t_{\mathrm{rip}} - t_0 = \frac{2 R_0^{-\epsilon/2}}{\epsilon B}.
\]
Identify $H_0 \equiv H(t_0) = B R_0^{\epsilon/2}$, so $R_0^{-\epsilon/2} = B/H_0$:
\[
t_{\mathrm{rip}} - t_0 = \frac{2}{\epsilon\,H_0}.
\]
Box the result:
\[
\boxed{\;t_{\mathrm{rip}} = t_0 + \frac{2}{\epsilon\,H_0}\quad(<\infty),\;}
\]
finite for any $\epsilon > 0$ --- the ``big rip'' singularity of phantom dark energy.

\end{document}
